\ncertchapter{4}{Projectiles}\label{ch:projectiles}

A projectile is any body moving under gravity alone once it has been launched: a thrown ball, a kicked
football, water from a hose, a long jumper in flight. It is the first real payoff of vectors, and it
turns on one observation that is easy to state and surprisingly hard to believe the first time: the
horizontal and vertical motions are completely independent of each other. A ball fired horizontally
from a table and a ball dropped from the same height at the same moment reach the floor together,
whatever the speed of the first. The first section explains why.

\begin{remarkbox}{The model and its limits}
Throughout this chapter air resistance is neglected, and $g$ is taken as constant in size and direction
(a flat Earth). This is an excellent model for a shot put, or for a slow ball over a few tens of metres. It
fails for a shuttlecock, a fast golf ball or an artillery shell: there drag shortens the range, makes
the descent steeper than the ascent, and lowers the best launch angle below $45^\circ$.
\end{remarkbox}

Once that is accepted, every projectile problem becomes two one-dimensional problems of the kind solved in
\chref{ch:straightline}, joined by the time. The chapter derives the familiar results for range, height and time of
flight, and then turns to the case those formulas do not cover: a launch from a height, where the
trajectory is no longer symmetric and the standard range formula gives the wrong answer.

\section{The Independence Principle}

Why should horizontal and vertical motion be independent at all? The answer rests on an experimental
fact: with air resistance negligible, every body moving freely near the Earth's surface has the same
acceleration, $g\approx\SI{9.8}{m/s^2}$ vertically downward (the free-fall acceleration of
\chref{ch:straightline}), whatever its speed or direction of motion.
Acceleration is the rate of change of velocity, and a vertical acceleration can change only the vertical
component of velocity; it has no horizontal part, so it cannot touch the horizontal motion. Each
component therefore obeys the constant-acceleration equations of \chref{ch:straightline} on its own:
horizontally with zero acceleration (uniform motion), vertically with acceleration $g$ downward, exactly
as if the body had been thrown straight up. The two share nothing but the clock. (Part B supplies the cause: after launch the only force is
gravity, which points straight down, and \chref{ch:newton} connects force to acceleration.)

\begin{keybox}{Horizontal and vertical motion are independent}
Once launched, a projectile has the constant acceleration $g$, vertically downward. Let it start with
velocity $\vv_0$ (a vector; the subscript $0$ marks a value at launch, $t=0$), of speed $v_0$ in
\si{m/s}, at angle $\theta$ above the horizontal, so that its components are
$v_{0x}=v_0\cos\theta$ and $v_{0y}=v_0\sin\theta$ (resolution, \chref{ch:vectors}). With the origin at the
launch point, $x$ the horizontal and $y$ the vertical (upward) coordinate, both in metres:
\[
\text{horizontal: } a_x=0,\quad x=v_{0x}\,t;\qquad
\text{vertical: } a_y=-g,\quad y=v_{0y}\,t-\tfrac12 g t^2,
\]
where $a_x$ and $a_y$ are the components of the acceleration; $a_y=-g$ because $y$ is measured upward.
Each is $s=ut+\tfrac12at^2$ of \chref{ch:straightline}, applied to one direction.
The two share only the clock $t$. A ball thrown horizontally and one dropped from the same height hit
the ground \emph{at the same instant}, because their \emph{vertical} motions are identical; the
horizontal velocity of the thrown ball does nothing to delay its fall.
\end{keybox}

The velocity follows from the same split. At time $t$ its components are $v_x=v_0\cos\theta$, which
never changes, and $v_y=v_0\sin\theta-gt$, which falls steadily; in Figure~\ref{fig:projectile} the
horizontal arrow is identical at every point drawn, while the vertical one shrinks to zero and reverses. The speed is
$v=\sqrt{v_x^2+v_y^2}$, and the velocity makes an angle $\tan^{-1}(v_y/v_x)$ with the horizontal. Because the velocity always points along the direction of motion, this arrow is
tangent to the path at every point: tilted upwards while the body rises, horizontal at the top, tilted
downwards on the way down.

\begin{trapbox}{the velocity at the highest point is not zero}
At the apex only the \emph{vertical} component vanishes. The body is still moving horizontally at
$v_0\cos\theta$, which is why it keeps advancing; only a ball thrown straight up ($\theta=90^\circ$)
is momentarily at rest. Nor is the acceleration zero there: it is $g$ downward at every instant of the
flight, apex included. If it were zero at the top, $v_y$ would stay zero and the ball would never come
down.
\end{trapbox}

\begin{figure}[htbp]
\centering
\tikzfig{kin_projectile}
\figcaption{fig:projectile}{The launch velocity $\vv_0$ splits into $v_0\cos\theta$ (horizontal) and $v_0\sin\theta$ (vertical). The horizontal component is the same arrow at launch, at the apex and on the way down; only the vertical component changes, and at the apex it is zero.}
\end{figure}

\section{Symmetric Launch: Range, Height, Time}

The independence principle gives $x(t)$ and $y(t)$ for any launch, but most questions ask for just
three numbers: how long the flight lasts, how high it goes and how far it travels. The simplest case,
and the one most often examined, is a launch from level ground. The trajectory is then
symmetric about its highest point, and the three quantities each take a compact form: the
\textbf{time of flight} $T$, the \textbf{maximum height} $H$ above the launch level, and the
\textbf{range} $R$, the horizontal distance from launch to landing.

\begin{derivbox}{Range, maximum height, time of flight}
Launch at speed $v_0$, angle $\theta$ above the horizontal, over level ground, with the origin at the
launch point.
\stepn{1} \textbf{Time of flight.} The height is $y=(v_0\sin\theta)t-\tfrac12gt^2=t\,(v_0\sin\theta-\tfrac12gt)$,
which is zero at launch and again when $T=\dfrac{2v_0\sin\theta}{g}$.
\stepn{2} \textbf{Range.} Horizontal distance in that time:
$R=v_0\cos\theta\cdot T=\dfrac{v_0^2\sin2\theta}{g}$ (using $2\sin\theta\cos\theta=\sin2\theta$).
\stepn{3} \textbf{Maximum height.} At the apex $v_y=0$, so from $v^2=u^2+2as$ applied vertically
($u=v_0\sin\theta$, $a=-g$, $s=H$), $0=(v_0\sin\theta)^2-2gH$:
$H=\dfrac{v_0^2\sin^2\theta}{2g}$.
\end{derivbox}

Those results give three points of the flight. To see the \emph{shape} of the whole path, and to read
launch data off a given curve, we need $y$ as a function of $x$ directly.

\begin{derivbox}{The trajectory is a parabola}
Eliminate $t$ between the two equations. From $x=(v_0\cos\theta)t$, $t=x/(v_0\cos\theta)$; substituting
in $y=(v_0\sin\theta)t-\tfrac12gt^2$ gives
\[
y=x\tan\theta-\frac{g x^2}{2v_0^2\cos^2\theta},
\]
which has the form $y=ax-bx^2$, with constants $a=\tan\theta$ and $b>0$ (here $a$ is just a
coefficient, not an acceleration): a downward-opening parabola (for $0\le\theta<90^\circ$). The
equation holds whether or not the ground is level, provided the origin is the launch point.
\end{derivbox}

\begin{keybox}{Consequences of $R=\dfrac{v_0^2\sin2\theta}{g}$}
For a given speed the range is greatest at $\theta=45^\circ$ (where $\sin2\theta=1$), and equals $v_0^2/g$.
Complementary angles give \emph{equal} range: $30^\circ$ and $60^\circ$ land at the same spot, because
$\sin(2\times30^\circ)=\sin(2\times60^\circ)$. Their \emph{heights}, however, differ sharply: since
$H\propto\sin^2\theta$ and $\sin30^\circ=\cos60^\circ$, the ratio is
$H_{60}/H_{30}=\sin^260^\circ/\cos^260^\circ=\tan^2 60^\circ=3$. The steep shot climbs three
times higher to reach the same target.
\end{keybox}

\begin{exbox}{A kicked football}
A football is kicked from level ground at \SI{25}{m/s} at $37^\circ$ above the horizontal (take
$\sin37^\circ=0.6$, $\cos37^\circ=0.8$, $g=\SI{10}{m/s^2}$). Find the time of flight, the maximum height,
the range, and the velocity at the apex and at $t=\SI{1}{s}$.

\stepn{1} \textbf{Resolve once.} $v_{0x}=25(0.8)=\SI{20}{m/s}$ and $v_{0y}=25(0.6)=\SI{15}{m/s}$.

\stepn{2} \textbf{Vertical motion alone.} $T=2v_{0y}/g=\SI{3}{s}$ and $H=v_{0y}^2/2g=225/20=\SI{11.25}{m}$.

\stepn{3} \textbf{Horizontal motion alone.} $R=v_{0x}T=20\times3=\SI{60}{m}$, in agreement with
$v_0^2\sin2\theta/g=625(0.96)/10$.

\stepn{4} \textbf{Velocity.} At the apex ($t=\SI{1.5}{s}$) $v_y=0$, so the velocity is \SI{20}{m/s}
horizontally, not zero. At $t=\SI{1}{s}$, $v_y=15-10=\SI{5}{m/s}$, so $v=\sqrt{20^2+5^2}\approx\SI{20.6}{m/s}$
at $\tan^{-1}(5/20)\approx14^\circ$ above the horizontal: still rising, but less steeply than at launch.
\ans{$T=\SI{3}{s}$; $H=\SI{11.25}{m}$; $R=\SI{60}{m}$; \SI{20}{m/s} at the apex}
\end{exbox}

\inwords{At any height $h$ between the launch level and the maximum, a projectile is there \emph{twice}:
once going up at time $t_1$, once coming down at $t_2$. For a launch from level ground $t_1+t_2=T$, the
total time of flight. The reason: $t_1$ and $t_2$ are the two roots of the quadratic
$\tfrac12gt^2-(v_0\sin\theta)t+h=0$, and the roots of $At^2+Bt+C=0$ sum to $-B/A$ and multiply to
$C/A$. So $t_1+t_2=2v_0\sin\theta/g$, and $t_1t_2=2h/g$, that is, $h=\tfrac12g\,t_1t_2$. This is the fastest route through many
exam problems: if you know the two times at which a projectile passes a given height, you know the flight
time without finding $v_0$ or $\theta$.}

\begin{trapbox}{the range formula needs level ground}
$R=v_0^2\sin2\theta/g$ assumes that launch and landing are at the \emph{same height}. For a launch from
a cliff, or a landing on a roof, go back to the independent $x$ and $y$ equations and solve for the
actual landing time.
\end{trapbox}

\section{Launch from a Height: The Asymmetric Case}

The trap just stated is common enough to deserve its own method. When launch and landing are at
\emph{different} heights, the neat symmetric formulas fail. If the landing
point is below the launch point, the projectile spends longer falling than it spent rising (and less
time if it lands above), so the path is not symmetric about the apex (compare the ascent and descent spans in
Figure~\ref{fig:projheight}). The reliable method is to write
the independent $x$ and $y$ equations and solve them.

\begin{figure}[htbp]
\centering
\tikzfig{kin_projectile_height}
\figcaption{fig:projheight}{A projectile launched upwards from height $H_0$. It rises to its apex, then falls \emph{past} the launch level to the ground. Because the horizontal speed never changes, the two spans marked along the ground are in the ratio of the times: the descent takes longer than the ascent, so the path is lopsided and the level-ground formulas for range and flight time do not apply.}
\end{figure}

\begin{methodbox}{A projectile launched from a height}
Take the launch point as origin, up positive, with speed $v_0$ at angle $\theta$; the launch point
is a height $H_0$ above the ground, so the ground is at $y=-H_0$.
\stepn{1} \textbf{Vertical (find the flight time).} Solve $-H_0=(v_0\sin\theta)\,t-\tfrac12 g t^2$ for
$t$. This is a quadratic; take the positive root. That $t$ is \emph{longer} than the level-ground
value $2v_0\sin\theta/g$, because of the extra drop $H_0$. For a horizontal throw ($\theta=0$) the
quadratic reduces to $H_0=\tfrac12gt^2$, so $t=\sqrt{2H_0/g}$, the same as for a body simply dropped.
\stepn{2} \textbf{Horizontal (find the range).} With that $t$, the horizontal distance is
$x=v_0\cos\theta\cdot t$.
\stepn{3} \textbf{Landing speed.} Use $v_x=v_0\cos\theta$ (unchanged) and $v_y=v_0\sin\theta-gt$, then
$v=\sqrt{v_x^2+v_y^2}$. Faster: the vertical motion falls a net $H_0$, so $v^2=u^2+2as$ with
$a=-g$ and $s=-H_0$ gives $v_y^2=(v_0\sin\theta)^2+2gH_0$, and then $v^2=v_x^2+v_y^2=v_0^2+2gH_0$, giving
$v=\sqrt{v_0^2+2gH_0}$, independent of the launch angle. (Energy conservation, in \chref{ch:energy},
gives the same result in one line.)
\end{methodbox}

\begin{exbox}{An oblique launch from a tower}
A ball is thrown at \SI{20}{m/s} at $30^\circ$ above the horizontal from the top of a tower \SI{15}{m}
high ($g=\SI{10}{m/s^2}$). Find the time of flight, the horizontal distance to the landing point, and the
landing speed.

\stepn{1} \textbf{Resolve.} $v_{0x}=20\cos30^\circ\approx\SI{17.3}{m/s}$, $v_{0y}=20\sin30^\circ=\SI{10}{m/s}$.

\stepn{2} \textbf{Vertical, origin at the launch point, up positive.} The ground is at $y=\SI{-15}{m}$:
\[
-15=10t-5t^2\quad\Longrightarrow\quad t^2-2t-3=0\quad\Longrightarrow\quad(t-3)(t+1)=0.
\]
The positive root is $t=\SI{3}{s}$. (The level-ground formula would give $2v_{0y}/g=\SI{2}{s}$: the extra
second is the fall below the launch level.)

\stepn{3} \textbf{Horizontal.} $x=v_{0x}t=17.3\times3\approx\SI{52}{m}$, against \SI{34.6}{m} from
$v_0^2\sin2\theta/g$, which would be wrong here.

\stepn{4} \textbf{Landing speed.} $v_x=\SI{17.3}{m/s}$ and $v_y=10-10(3)=\SI{-20}{m/s}$, so
$v=\sqrt{17.3^2+20^2}\approx\sqrt{700}\approx\SI{26.5}{m/s}$, in agreement with
$\sqrt{v_0^2+2gH_0}=\sqrt{400+300}$. The ball strikes the ground at $\tan^{-1}(20/17.3)\approx49^\circ$
below the horizontal.
\ans{$t=\SI{3}{s}$; $x\approx\SI{52}{m}$; $v\approx\SI{26.5}{m/s}$}
\end{exbox}

\section{Summary}\label{sec:summary4}
Once launched, a projectile feels only gravity, so its horizontal velocity is constant and its vertical
motion is uniformly accelerated at $g$. The two are independent and share only the time. For a launch
from level ground the trajectory is a symmetric parabola, and the three standard results follow directly.
For a launch from a height none of the symmetric formulas applies, and the safe method is to write
$x(t)$ and $y(t)$ separately and solve.

\begin{itemize}
\item $x=v_0\cos\theta\,t$, \quad $y=v_0\sin\theta\,t-\tfrac12gt^2$.
\item Level ground: $T=\dfrac{2v_0\sin\theta}{g}$, \quad $H=\dfrac{v_0^2\sin^2\theta}{2g}$, \quad
      $R=\dfrac{v_0^2\sin2\theta}{g}$.
\item The range is greatest at $45^\circ$ for a given speed, and complementary angles ($30^\circ$ and $60^\circ$) give equal
      ranges but very different heights and flight times.
\item The trajectory equation $y=x\tan\theta-\dfrac{gx^2}{2v_0^2\cos^2\theta}$ lets the launch parameters be
      read off a given parabola.
\item At the highest point the velocity is horizontal, not zero; the acceleration is still $g$ downward.
\item The range formula assumes level ground. From a height, solve the quadratic in $t$ and keep the
      positive root.
\end{itemize}

\section{Problem Types}\label{sec:types4}
Five problem types account for almost every projectile question. All five yield to the same discipline:
write the horizontal and vertical equations separately, use the time to connect them, and resist the
temptation to apply a level-ground formula to a launch that is not from level ground.

\archq{Find $T$, $H$ or $R$ from the launch data}
{``A body is projected at speed $v_0$ and angle $\theta$; find the time of flight, maximum height or
range.''}
{Resolve once: $v_{0x}=v_0\cos\theta$ (constant), $v_{0y}=v_0\sin\theta$ (decelerated by $g$). Then
$T=\dfrac{2v_0\sin\theta}{g}$, $H=\dfrac{v_0^2\sin^2\theta}{2g}$, $R=\dfrac{v_0^2\sin2\theta}{g}$. For ratio
questions never find $v_0$; divide the two expressions.}
{At the top the velocity is \emph{not} zero: only the vertical component vanishes; the horizontal
component $v_0\cos\theta$ survives untouched. Also remember $\theta$ and $(90^\circ-\theta)$ give the
\emph{same} range.}
{A ball is projected at \SI{20}{m/s} at $30^\circ$ to the horizontal ($g=\SI{10}{m/s^2}$). Its time of flight and
maximum height are:}
{\SI{2}{s}, \SI{5}{m}}{\SI{1}{s}, \SI{5}{m}}{\SI{2}{s}, \SI{10}{m}}{\SI{4}{s}, \SI{20}{m}\archans{A}{
$v_{0y}=20\sin30^\circ=\SI{10}{m/s}$. $T=2v_{0y}/g=\SI{2}{s}$; $H=v_{0y}^2/2g=100/20=\SI{5}{m}$.}}

\archq{Horizontal projection from a height}
{``A body is thrown horizontally at $v_0$ from height $h$; find the time to land, the horizontal
distance, or the striking speed.''}
{The vertical motion is pure free fall, so $t=\sqrt{2h/g}$ \emph{independent of $v_0$}. Then
range $=v_0t$, and the landing speed is $\sqrt{v_0^2+(gt)^2}$.}
{Thinking that a faster horizontal throw stays in the air longer. It does not: neglecting air
resistance, a ball thrown horizontally and one dropped from the same height land at the \emph{same
instant}.}
{A stone is thrown horizontally at \SI{15}{m/s} from a \SI{45}{m} cliff ($g=\SI{10}{m/s^2}$). How far from the base
does it land?}
{\SI{30}{m}}{\SI{60}{m}}{\SI{90}{m}}{\SI{45}{m}\archans{D}{$t=\sqrt{2(45)/10}=\SI{3}{s}$;
horizontal distance $=15\times3=\SI{45}{m}$.}}

\archq{Speed or velocity at a stated point of the flight}
{``Find the speed at the highest point'', or ``the velocity after $t$ seconds'', or ``the speed on
striking the ground''.}
{Track the two components separately. $v_x=v_0\cos\theta$ is \textbf{constant throughout};
$v_y=v_0\sin\theta-gt$ changes. Then $v=\sqrt{v_x^2+v_y^2}$. At the apex $v_y=0$, so the speed there is
$v_0\cos\theta$, \emph{not} zero.}
{Saying the velocity at the top is zero. Only the \emph{vertical} component vanishes; the projectile is
still moving horizontally, which is why it keeps advancing.}
{A body is projected at \SI{20}{m/s} at $60^\circ$ to the horizontal. Its speed at the highest point is:}
{zero}{\SI{17.3}{m/s}}{\SI{10}{m/s}}{\SI{20}{m/s}\archans{C}{At the apex only $v_x$ survives:
$v_x=20\cos60^\circ=20(0.5)=\SI{10}{m/s}$. Option (B) is $20\sin60^\circ$, the wrong component.}}

\archq{Equal ranges at complementary angles}
{``Two projectiles with the same speed have the same range'', or ``compare the heights or times for
$\theta$ and $90^\circ-\theta$''.}
{Since $\sin2\theta=\sin\big(2(90^\circ-\theta)\big)$, the angles $\theta$ and $(90^\circ-\theta)$ give
\textbf{equal ranges}. But the heights and times are \emph{not} equal: with $H_1,T_1$ for $\theta$ and
$H_2,T_2$ for $90^\circ-\theta$, $\dfrac{H_1}{H_2}=\tan^2\theta$ and $\dfrac{T_1}{T_2}=\tan\theta$.}
{Assuming equal range means equal everything. The steeper launch always gives the greater height and the
longer flight.}
{Two balls are projected at the same speed at $30^\circ$ and $60^\circ$. The ratio of their maximum
heights $H_{30}:H_{60}$ is:}
{$1:1$}{$1:3$}{$3:1$}{$1:\sqrt3$\archans{B}{$H\propto\sin^2\theta$, so the ratio is
$\sin^230^\circ:\sin^260^\circ=(1/4):(3/4)=1:3$. Their ranges, however, are equal.}}

\archq{Read the parameters off the trajectory equation}
{``The path of a projectile is $y=ax-bx^2$; find the range, the maximum height or the launch angle.''}
{Compare with the standard form $y=x\tan\theta-\dfrac{gx^2}{2v_0^2\cos^2\theta}$, so $a=\tan\theta$ and
$b=\dfrac{g}{2v_0^2\cos^2\theta}$. The range comes from setting $y=0$: $R=a/b$. The maximum height is
where $dy/dx=a-2bx=0$, at $x=a/2b=R/2$, and there $y=H=\dfrac{a^2}{4b}$.}
{Trying to reconstruct $v_0$ and $\theta$ before answering. For range and height you never need them;
read the answers straight off $a$ and $b$.}
{A projectile follows $y=x-\dfrac{x^2}{20}$ (metres). Its horizontal range is:}
{\SI{20}{m}}{\SI{10}{m}}{\SI{40}{m}}{\SI{5}{m}\archans{A}{Setting $y=0$:
$x\left(1-\dfrac{x}{20}\right)=0$, giving $x=0$ (launch) and $x=\SI{20}{m}$ (landing).}}
